\section{Conclusion}

We introduce TPE, a three-stage temporal uncertainty localization
framework that closes the long-form / real-time gap for mechanistic
hallucination detection. The encoder reads the full hidden-state
activations; the temporal head produces per-token logits under weakly
supervised top-$k$ MIL within each claim; the passage-level
refinement step turns those logits into per-claim scores that
discriminate uncertain claims from confident ones within the same
response. On RAGTruth, FactScore, and LongFact across three
instruction-tuned LLMs, the per-token output reaches token- and
sentence-level AUROC competitive with pooled-activation baselines while
running $25\times$ faster on long passages. On the
LM-Polygraph short-form benchmark, the same model produces a
sequence-level score plug-compatible with existing baselines. Promising
directions include richer attention-side features and shared per-LLM
adapters across model families.

\paragraph{Limitations.} The auxiliary span labels we train on are
noisy: claim-level verification pipelines mis-attribute hallucinated tokens to
neighbouring supported tokens, and the relevant activation signal can
fall outside the span (Appendix Figure~\ref{fig:signal_propagation}). The
Stage~3 BCE absorbs some of this noise but the per-claim numbers
still trail the encoder-only Stage~3 variant on the window ablation,
suggesting the temporal head over-fits to span boundaries when the
supervision is locally inconsistent. Our LLM-specific Linear Adapter
also assumes white-box access to hidden states and a small calibration
set per LLM; closed-weight settings remain out of scope. Finally, we
evaluate on English corpora and instruction-tuned chat models;
generalization to other languages and to base or reasoning-tuned
checkpoints is left for future work.
